\documentclass{article}

\usepackage[eandd]{neurips_2026}

\usepackage[utf8]{inputenc}
\usepackage[T1]{fontenc}
\usepackage{hyperref}
\usepackage{url}
\usepackage{booktabs}
\usepackage{amsfonts}
\usepackage{amsmath}
\usepackage{amssymb}
\usepackage{nicefrac}
\usepackage{microtype}
\usepackage{xcolor}
\usepackage{multirow}
\usepackage{graphicx}
\usepackage{siunitx}

\DeclareMathOperator*{\argmax}{arg\,max}

\newcommand{\dataset}{\textsc{CTSpinoPelvic1K}}
\newcommand{\TS}{TotalSegmentator}

\title{Diagnosing the LSTV Anomaly Gap: A Stratified
Evaluation Protocol Exposing Failure Modes in Production CT
Segmentation Systems}

% Anonymized for double-blind review (default for E&D Track)
\author{%
  Anonymous Author(s) \\
  Anonymous Affiliation \\
  \texttt{anonymous@institution.edu} \\
}

\begin{document}

\maketitle

% =============================================================================
\begin{abstract}
\TS{} (TS) is the most widely deployed multi-organ CT segmentation
model in clinical and research practice and reports a headline
test-set Dice of 0.943.
We show that on Lumbosacral Transitional Vertebrae (LSTV)---a
congenital variant present in 5--35\% of the population and the
single highest-risk anatomic feature for wrong-level spine
surgery---TS junction Dice drops by \textbf{29.0 points}
(0.824\,$\to$\,0.534) on Any-LSTV cases relative to normal
anatomy, and by \textbf{46.8 points} (0.824\,$\to$\,0.357)
specifically on lumbarization.
We characterize a previously undocumented mechanistic failure mode:
\emph{vertebral level-shift mislabeling}, in which TS, lacking an
L6 output class, distributes 5 labels across 6 lumbar bodies by
shifting label assignment by one anatomic level.
On lumbarization cases, GT L6 voxels are predominantly labeled as
L5, GT L5 as L4, and so on, producing per-class DSCs that
collapse caudally: L1=0.878 → L4=0.388 → L5=0.291.
On count-style sacralization (4 lumbar bodies + fused L5), the
inverse failure occurs: TS predicts L1--L5 in standard positions
while GT spans only L1--L4, collapsing per-class DSC to near
zero.
TS structurally cannot predict L6 (\textsc{TS\_L6}\,=\,0 across
all 1{,}152 evaluation cases), making this an architectural
failure compounded by distributional underrepresentation.
This failure is invisible to every existing CT segmentation
benchmark because none stratify evaluation by transitional
anatomy.
We make three contributions to close this gap:
(i) \textbf{\dataset{}}, a CT-native benchmark of 1{,}153 abdominal
CT volumes across 802 patients with unified 10-class spinopelvic
masks and 33 LSTV-positive cases spanning a 6-way subtype
taxonomy with explicit dual-origin arbitration metadata;
(ii) a \textbf{stratified evaluation protocol} reporting per-class
DSC, junction-DSC over a 40\,mm L5/S1 axial window, and per-class
confusion matrices, designed to make architectural failures
(no L6 output class) attributable separately from distributional
failures (LSTV underrepresentation in training);
(iii) a \textbf{full-resolution H200 baseline} demonstrating
the gap is closable: a 5-fold CV nnU-Net trained at native
$0.75\!\times\!0.75\!\times\!0.80$~mm with a
$256\!\times\!320\!\times\!320$ patch (the only configuration
fitting the full lumbar spine and pelvis on a single GPU,
requiring 100~GB HBM3e) recovers junction Dice on LSTV cases to
0.\textbf{XXX} (\textbf{+XX.X} points) and emits an L6 class on
all lumbarization cases.
The hardware budget is not incidental: TS's published
1.5~mm and 3~mm resolutions, which fit on commodity GPUs,
discard the boundary detail that distinguishes Castellvi
Type II morphology from normal L5/S1.
\dataset{}, the evaluation protocol, and trained checkpoints
will be released publicly upon acceptance.
\end{abstract}

% =============================================================================
\section{Introduction}
\label{sec:intro}

Modern radiology depends on automated segmentation.
\TS{}~\citep{Wasserthal2023TotalSegmentator:Images}, an nnU-Net-based
model trained on 1{,}204 routine CT examinations to delineate
104 anatomic structures, has become the de-facto production
system for whole-body CT segmentation, with thousands of clinical
and research installations and integration into commercial
radiology workstations.
The original paper reports a test-set Dice of 0.943 across all
classes and qualitatively documents two known failure modes for
the spine: ``mixing neighboring vertebrae classes'' (1\% of cases)
and ``mixing neighboring rib classes''
(12\%)~\citep[Figure 4]{Wasserthal2023TotalSegmentator:Images}.
For most anatomy, this is acceptable headline performance.

This paper concerns one anatomic class on which it is not.

The Lumbosacral Transitional Vertebra (LSTV) is a congenital
variant in which the L5/S1 junction exhibits morphological features
of the adjacent segment~\citep{Castellvi1984LumbosacralDefects}: \emph{sacralization},
where L5 fuses with the sacrum, or \emph{lumbarization}, where S1
detaches and presents as a discrete sixth lumbar body (L6).
Population prevalence ranges from 5\% to
35\%~\citep{Sekharappa2014LumbosacralSignificance} depending on imaging modality
and classification criterion, and LSTV misidentification is the
single largest contributor to wrong-level spine
surgery~\citep{Konin2010LumbosacralRelevance,Carrino2011EffectImaging}.
A model that fails on LSTV is not failing on a benign edge case;
it is failing on the exact anatomy whose accurate delineation
defines its clinical value.

\paragraph{The evaluation gap.}
TS does not report LSTV-stratified performance.
Neither does VerSe~\citep{Sekuboyina2020VerSe:Benchmark},
CTSpine1K~\citep{Deng2021CTSpine1K:Tomography}, or
CTPelvic1K~\citep{Liu2021DeepModels}, the three other major CT spine
and pelvis benchmarks.
This is not an oversight in any single paper:
\emph{no public CT segmentation benchmark stratifies evaluation
by transitional anatomy}, because no public CT segmentation
benchmark has curated LSTV cases at the scale required to power
such analysis.
The result is a measurement vacuum: we cannot know how badly TS
or any successor fails on LSTV until a stratified benchmark
exists.

\paragraph{The mechanistic failure mode.}
We document a specific failure mode in TS on LSTV cases that the
original paper does not characterize:
\textbf{vertebral level-shift mislabeling}.
TS's label vocabulary contains
\texttt{vertebrae\_L1}--\texttt{vertebrae\_L5} but no L6 class.
On lumbarization (a 6-vertebra spine), TS does not output
``unknown'' on the supernumerary body; instead, it shifts the
entire lumbar label assignment by one anatomic level so that all
five outputs remain in use.
Empirically this manifests as: GT L6 voxels predicted predominantly
as L5; GT L5 voxels predicted predominantly as L4; cascading
upward.
Per-class DSC consequently collapses caudally
(L1=0.878\,$\to$\,L4=0.388\,$\to$\,L5=0.291)
while remaining near-normal for the most rostral lumbar bodies
where the topological constraint is least binding.
The output is locally well-formed---each predicted vertebra is
spatially compact, anatomically positioned, indistinguishable
from a single correct prediction except by counting and by
boundary alignment with the actual lumbar bodies.
A counting-aware downstream system---a surgical planning pipeline
that reads ``L4--L5 disc'' from segmentation output, for
instance---will fetch the wrong disc.

The same architectural mismatch produces an inverse failure on
count-style sacralization (4 visible lumbar bodies plus a fused
L5/sacrum mass).
GT here ranges only L1--L4; TS, expecting 5 lumbar bodies,
emits L1--L5 in standard positions, producing systematic
$+1$-level mislabeling from L2 onward.
Per-class DSC collapses on every lumbar class
(L1=0.265, L2=0.100, L3=0.055, L4=0.001, L5=0.000).
The failure is mechanically attributable to a single architectural
decision (no L6 output class) compounded by a learned response to
distributional mismatch (LSTV underrepresentation in training);
it is invisible to every metric except a stratified per-class DSC
plus per-class confusion analysis at the L5/S1 junction.

\paragraph{Closing the gap requires hardware.}
Demonstrating that the gap is closable requires a baseline that
both incorporates LSTV cases in training \emph{and} fits sufficient
anatomical context per training patch to learn the spinopelvic
spatial relationships at the L5/S1 junction.
At spacing $0.75\!\times\!0.75\!\times\!0.80$~mm and patch size
$256\!\times\!320\!\times\!320$ (the only patch size fitting both the
pelvis and the full lumbar spine in one window on the resampled
COLONOG geometry), a single training batch consumes
${\sim}100$~GB of GPU memory under nnU-Net's
\texttt{nnUNetResEncUNetPlans\_100G} planner.
This budget is available on H200 hardware (140~GB HBM3e) but not
on the L40S/A100/H100 GPUs that bracket TS's reported training
configuration and most academic baselines.
The hardware budget is the difference between a model that sees
the L5--sacrum--ilium relationship during training and one that
does not.
TS's published training used patches at 1.5~mm and 3~mm isotropic
resolution, which fit comfortably on commodity GPUs but discard
the sub-millimeter boundary detail required to disambiguate
Castellvi Type~II from a normal L5/S1.

\paragraph{Contributions.}
\begin{enumerate}
\item \textbf{LSTV failure-mode analysis of \TS{} on CT.}
  The first stratified evaluation of TS on LSTV-positive cases,
  with per-class DSC, junction-DSC, and per-class confusion
  matrices reported across a 6-way LSTV taxonomy.
  TS junction-DSC drops 29.0 points on Any-LSTV cases and 46.8
  points specifically on lumbarization. Confusion matrices reveal
  the level-shift mechanism with case-level documentation
  (\S\ref{sec:case_studies}) showing GT-L6 voxels labeled as L5
  in 81{,}088 / 85{,}186 voxels (95\%) on a representative case.
  \textsc{TS\_L6}\,=\,0 across all 1{,}152 evaluation cases,
  confirming the architectural component.
  TS is evaluated in full-resolution mode, isolating the failure
  as a property of label vocabulary and training distribution
  rather than inference resolution.
\item \textbf{\dataset{}: a stratified CT spinopelvic benchmark.}
  1{,}153 abdominal CT records across 802 matched patients,
  unified 10-class spinopelvic masks (with an explicit L6 class),
  33 LSTV-positive cases spanning a 6-way subtype taxonomy
  (lumbarization $n=14$; sacralization-by-count $n=9$;
  semi-sacralization $n=2$; full sacralization $n=6$;
  ambiguous $n=2$), dual-origin LSTV labels with
  explicit arbitration metadata for the systematically-ambiguous
  lumbosacral confusion zone, and a registration-free fusion
  pipeline that resolves a previously undocumented Y-axis
  coordinate-convention mismatch between CTSpine1K and CTPelvic1K.
\item \textbf{Stratified evaluation protocol.}
  Per-LSTV-subgroup DSC, junction-DSC over a 40~mm axial window
  centered at L5/S1, and per-class voxel confusion matrices.
  Together these make architectural and distributional failure
  modes individually attributable.
\item \textbf{Full-resolution H200 baseline closing the gap.}
  An nnU-Net 3d\_fullres baseline trained at native
  $0.75\!\times\!0.75\!\times\!0.80$~mm spacing with a
  $256\!\times\!320\!\times\!320$ patch under the
  \texttt{nnUNetResEncUNetPlans\_100G} planner targeting 100~GB
  GPU memory, on H200 NVL hardware, with 5-fold CV across all
  979 training cases via partial-annotation training.
  The baseline closes the LSTV junction-Dice gap to
  $\Delta = $\textbf{XX} points and emits an L6 class on all
  lumbarization cases in the held-out test set.
\end{enumerate}

% =============================================================================
\section{Related Work}
\label{sec:related}

\paragraph{TS{} and the production-baseline problem.}
\citet{Wasserthal2023TotalSegmentator:Images} segments 104 anatomic
structures from CT and is widely deployed.
The original paper documents ``mixing neighboring vertebrae
classes'' (1\% of cases) as a failure mode but does not stratify
by transitional anatomy or characterize level-shift labeling
behavior. TS's failure modes on LSTV-positive CT have, to our
knowledge, not been previously characterized.
\citet{Warszawer2025TotalSpineSeg:MRI} extended this methodology
to MRI with spine-specific architecture, but MRI cannot resolve
the cortical bridging that defines Castellvi morphology.

\paragraph{CT spine and pelvis benchmarks.}
CTSpine1K~\citep{Deng2021CTSpine1K:Tomography} provides 1{,}005 CT volumes
with per-vertebra annotation;
VerSe~\citep{Sekuboyina2020VerSe:Benchmark} provides 460 challenge-format
spine CTs; CTPelvic1K~\citep{Liu2021DeepModels} provides 1{,}184
pelvic CTs.
None reports LSTV prevalence or stratifies splits by transitional
anatomy.
CTPelvic1K filename qualifiers qualitatively encode an LSTV
phenotype for many cases, which we systematically extract,
disambiguate (separating semi-sacralization from full
sacralization, a distinction lost in earlier extraction
attempts due to a substring-match parser bug), and validate
against CTSpine1K vertebral counting.

\paragraph{Multi-source label fusion.}
\citet{Li2025LumbarMechanism} combine CTSpine1K and CTPelvic1K for
joint segmentation but treat them as two independent sources
without resolving the silent coordinate-convention mismatch we
describe in \S\ref{sec:dataset}, and provide no LSTV-stratified
evaluation.

\paragraph{Evaluation as scientific object.}
The position that benchmarks should report stratified performance
on safety-critical subgroups, rather than relying on aggregate
metrics, is established in the medical ML
community~\citep{Sang2025BenchmarkChallenge,Pieper2024Spine1}.
This paper applies that principle to a setting where the relevant
subgroup (LSTV) is invisible to every existing benchmark, and
demonstrates that adding the stratification exposes a quantifiable
failure mode in the most widely deployed production model.

% =============================================================================
\section{Failure-Mode Analysis: \TS{} on LSTV}
\label{sec:ts_failure}

We evaluate TS v2 (\texttt{total} task) zero-shot at
\textbf{full resolution} (\texttt{--fast} disabled) on the full
\dataset{} matched cohort.
TS was not trained on any \dataset{} case, so every case is a
valid test case; this provides $N=1{,}152$ successful evaluation
records aggregated to 801 patients via patient-level merging of
separate-mode (spine + pelvic) records, including 33 LSTV-positive
patients across the 6-way taxonomy.
Full-resolution mode isolates the failure as a property of label
vocabulary and training distribution rather than inference
resolution, addressing the central confound of low-resolution
benchmarks.

\subsection{Stratified Per-Class Performance}

Table~\ref{tab:ts_stratified} reports TS performance stratified by
LSTV phenotype using the 6-way taxonomy.
On normal anatomy ($n=768$) TS performs at its published level:
mean L1--L5 DSC of 0.929, sacrum 0.817, hips 0.928, junction-DSC
0.824.
On Any-LSTV ($n=33$), every metric drops:
junction-DSC falls $\Delta = -29.0$ points to 0.534.
The drop is highly asymmetric across phenotypes.
Lumbarization ($n=14$) collapses caudally (L4=0.388, L5=0.291)
with junction-DSC at 0.357.
Sacralization-by-count ($n=9$) collapses across all lumbar classes
(L1=0.265 down to L5=0.000) but retains junction-DSC at 0.697
because the surviving GT L1--L4 voxels are sparsely distributed
within the 40~mm window.
Full sacralization ($n=6$) and semi-sacralization ($n=2$) retain
moderate-to-high per-class DSC, indicating that morphological-
qualifier-based phenotypes do not propagate the same level-shift
mislabeling pattern.

\begin{table}[t]
\caption{TS{} zero-shot stratified by LSTV phenotype on
\dataset{}, full-resolution mode.
Junction DSC computed over a 40~mm axial window centered at
L5/S1.
\textsc{TS\_L6} = number of cases in which TS emitted any L6
voxel (structurally zero by label vocabulary).
JxnDSC entries marked with $\dagger$ indicate fused-mode patient
unavailable; semi-sacralization and ambiguous have only
separate-mode records, where junction analysis is not computed.
The junction-DSC drop on lumbarization is mechanistically driven
by the level-shift labeling pattern characterized in
\S\ref{sec:level_shift}, not by uniform performance degradation
across classes.}
\label{tab:ts_stratified}
\centering
\small
\begin{tabular}{lrcccccccr}
\toprule
\bfseries Subgroup & \bfseries $n$ &
\bfseries L1 & \bfseries L2 & \bfseries L3 & \bfseries L4 & \bfseries L5 &
\bfseries Sac & \bfseries Hips &
\bfseries Jxn DSC \\
\midrule
All cases                              & 801 & 0.911 & 0.906 & 0.913 & 0.914 & 0.913 & 0.817 & 0.929 & 0.818 \\
Normal                                 & 768 & 0.921 & 0.918 & 0.929 & 0.936 & 0.938 & 0.817 & 0.929 & 0.824 \\
\midrule
Any LSTV                               & 33  & 0.681 & 0.615 & 0.523 & 0.367 & 0.314 & 0.822 & 0.929 & 0.534 \\
\quad Lumbarization                    & 14  & 0.878 & 0.853 & 0.690 & 0.388 & 0.291 & 0.833 & 0.928 & \textbf{0.357} \\
\quad Sacralization (count-style)      &  9  & 0.265 & 0.100 & 0.055 & 0.001 & 0.000 & 0.814 & 0.926 & 0.697 \\
\quad Semi-sacralization               &  2  & 0.951 & 0.951 & 0.951 & 0.955 & 0.956 & 0.824 & 0.937 & $\dagger$ \\
\quad Sacralization (full)             &  6  & 0.718 & 0.676 & 0.673 & 0.724 & 0.793 & 0.826 & 0.932 & 0.636 \\
\quad Ambiguous                        &  2  & 0.952 & 0.951 & 0.882 & 0.678 & 0.374 & 0.771 & 0.932 & $\dagger$ \\
\midrule
\textsc{TS\_L6} (any subgroup)         &     &       &       &       &       &       &       &       & 0 / 1152 \\
\bottomrule
\end{tabular}
\end{table}

\subsection{The Level-Shift Mislabeling Failure}
\label{sec:level_shift}

The asymmetry between lumbarization (severe, caudally-progressive
collapse) and sacralization-by-count (uniform collapse) is
mechanistically explainable by a single architectural fact: TS
has no L6 output class.

For lumbarization (6 lumbar bodies), the fixed label cardinality
forces TS to compress 6 anatomic units into 5 labels.
The most rostral body remains correctly labeled L1 because the
upward escape (into thoracic vertebrae T11--T12, which TS does
have) preserves the local label-to-anatomy correspondence.
But every successive body is progressively misaligned, with the
final body labeled L5 (correctly the position for the 5th lumbar
body) when in fact it occupies the L6 position.
The 5th-from-top body, which is anatomically L5, gets labeled L4.
Hence the caudally-progressive DSC collapse:
L1=0.878\,$\to$\,L2=0.853\,$\to$\,L3=0.690\,$\to$\,L4=0.388\,$\to$\,L5=0.291.
The L1 retains 0.878 (a 0.04 drop from normal) only because the
top body is genuinely L1; L4 and L5 collapse because the labels
they receive belong on the bodies one position higher.

For count-style sacralization (4 lumbar bodies + fused L5),
the failure is the inverse.
Ground truth here labels only L1--L4, with the fused L5 absorbed
into the sacrum class.
TS, expecting 5 lumbar bodies, emits L1--L5 in standard positions.
The five predicted labels overlap the four GT lumbar bodies with a
$+1$-level offset (predicted L2 sits where GT L1 is, etc.), so
DSC collapses uniformly across L1--L5
(0.265, 0.100, 0.055, 0.001, 0.000).
Junction-DSC remains at 0.697 because sacrum and the surviving
caudal lumbar voxels still align reasonably within the 40~mm
window---these cases lose lumbar identity but not sacral
boundaries.

Full sacralization ($n=6$, retains junction-DSC 0.636 and L1--L5
DSCs 0.67--0.79) shows neither pattern.
We interpret this as a less-systematic failure mode:
full-sacralization GT typically labels the fused L5 mass partly
as L5 and partly as sacrum, giving TS's L5 prediction enough
overlap that DSC remains substantial.
Semi-sacralization ($n=2$, L1--L5 all $>0.95$) shows essentially
no failure; the morphological subtlety distinguishing semi from
normal does not impose the topological constraint that breaks
TS's label vocabulary.

\subsection{Mechanistic versus Distributional Attribution}

We attribute the failure to two compounding causes that the
stratified protocol disentangles.

\textbf{(i) Architectural:} TS's label vocabulary excludes L6, so
any model with TS's output head cannot correctly segment a
lumbarized spine regardless of training data; \textsc{TS\_L6} is
structurally zero across all 1{,}152 evaluation cases.
This contributes a fixed lower bound on lumbarization DSC.

\textbf{(ii) Distributional:} TS's training data
(1{,}204 hospital CTs sampled from 2012, 2016, 2020 with no LSTV
stratification) likely contains a population-base-rate fraction of
LSTV cases, but these are too few and too unstratified to teach
the model alternative compensatory behaviors (e.g., explicit
``unknown'' or low-confidence prediction for the supernumerary
body).
The level-shift behavior is a learned response to a
distributional mismatch (more lumbar bodies than the trained
distribution), constrained by the architectural bound.

A baseline trained with an L6 class on LSTV-stratified data
should close the architectural component completely
(\textsc{baseline\_L6}\,$>0$ on lumbarization cases).
Whether the level-shift behavior also resolves depends on whether
LSTV oversampling provides sufficient signal at the relevant
resolution; we test both in \S\ref{sec:baseline}.

% =============================================================================
\section{Case Studies: Failure Modes of \TS{} on LSTV}
\label{sec:case_studies}

The aggregate metrics in Table~\ref{tab:ts_stratified} summarize a
population-level pattern but obscure case-level mechanism.
This section presents three case studies---one each for severe
lumbarization, mild lumbarization with adaptive sacrum
labeling, and count-style sacralization---using per-class voxel
confusion matrices over the 40~mm L5/S1 window to expose the
exact mechanism of failure.
Each case is a real \dataset{} record evaluated by TS under the
same configuration as Table~\ref{tab:ts_stratified}.

\subsection{Case 1: Lumbarization with Severe Level Shift
            (Token 167)}
\label{sec:case167}

Token 167 is a fused-mode lumbarization case with junction-DSC
0.221, the lowest among lumbarization cases.
The per-class confusion matrix over the 40~mm L5/S1 window reveals
the level-shift mechanism cleanly
(Table~\ref{tab:case167_confusion}).
Of 92{,}629 GT-L5 junction voxels, 92{,}537 (99.9\%) are predicted
as L4---essentially the entire L5 region is misnamed.
Of 85{,}186 GT-L6 voxels, 81{,}088 (95.2\%) are predicted as L5
and only 9 as L4; the remaining 4{,}089 fall to background.
The sacrum prediction is largely intact: 80{,}334 of 149{,}005
GT-sacrum voxels (53.9\%) are correctly labeled, with the
remainder going to background due to incomplete sacrum-extent
prediction (a separate, milder issue).
No GT-L6 voxel is predicted as L6 (\textsc{TS\_L6}\,=\,0
structurally).

\begin{table}[t]
\caption{Case 1 (Token 167) per-class voxel confusion matrix over
the 40~mm L5/S1 axial window. Rows = GT class; columns = TS
prediction. Bold cells indicate the dominant prediction class for
each GT row. The diagonal-shifted dominance pattern (GT-L5
$\to$ predicted L4; GT-L6 $\to$ predicted L5) is the empirical
signature of the level-shift mislabeling failure.}
\label{tab:case167_confusion}
\centering
\small
\begin{tabular}{lrrrrrr}
\toprule
\bfseries GT $\backslash$ Pred & \bfseries bg & \bfseries L3 & \bfseries L4 & \bfseries L5 & \bfseries sacrum & \bfseries L6 \\
\midrule
L5     & 4{,}572  & 2 & \textbf{92{,}537} & 90              & 0           & 0 (n/a) \\
L6     & 4{,}089  & 0 & 9                 & \textbf{81{,}088} & 0         & 0 (n/a) \\
sacrum & 68{,}635 & 0 & 0                 & 36              & \textbf{80{,}334} & 0 (n/a) \\
\bottomrule
\end{tabular}
\end{table}

The output is anatomically well-formed: TS produces five compact
vertebral predictions in roughly correct positions plus a sacrum
prediction.
A reader unaware of the patient's lumbarization could not detect
the failure from a single sagittal slice---each labeled vertebra
``looks like'' the labeled vertebra.
The failure is detectable only by counting bodies (six lumbar,
not five) or by comparing predicted-label-to-position correspondence
to the patient's actual anatomy.
A surgical planner consuming this output and instructed to operate
at L4--L5 would target the disc between predicted-L4 and
predicted-L5, which is anatomically the L5--L6 disc---one level
caudal to the intended target.

\subsection{Case 2: Lumbarization with L6 Suppression
            (Token 189)}
\label{sec:case189}

Token 189 is a fused-mode lumbarization case with the highest
junction-DSC among lumbarization cases (0.617), demonstrating a
qualitatively different failure mode.
The confusion matrix (Table~\ref{tab:case189_confusion}) shows
that GT-L5 voxels are predominantly predicted correctly as L5
(135{,}224 of 140{,}564 = 96.2\%), and GT-sacrum voxels are
correctly predicted as sacrum (227{,}881 of 247{,}665 = 92.0\%).
The L6 region, however, is predicted predominantly as
\emph{background} (110{,}322 of 111{,}448 = 99.0\%).

\begin{table}[t]
\caption{Case 2 (Token 189) per-class voxel confusion matrix over
the 40~mm L5/S1 axial window. The L6 row shows the L6-suppression
failure mode: TS labels the supernumerary body as background
rather than committing it to a wrong lumbar level.}
\label{tab:case189_confusion}
\centering
\small
\begin{tabular}{lrrrrrr}
\toprule
\bfseries GT $\backslash$ Pred & \bfseries bg & \bfseries L4 & \bfseries L5 & \bfseries sacrum & \bfseries L6 \\
\midrule
L5     & 4{,}235          & 1{,}105 & \textbf{135{,}224} & 0                  & 0 (n/a) \\
L6     & \textbf{110{,}322} & 0     & 907                & 219                & 0 (n/a) \\
sacrum & 19{,}683         & 0       & 101                & \textbf{227{,}881} & 0 (n/a) \\
\bottomrule
\end{tabular}
\end{table}

This is functionally a different mistake from Case 1.
Rather than shifting labels by one anatomic level, TS appears
to recognize that an unfamiliar body is present and declines to
label it---a kind of high-uncertainty refusal.
The L5 and sacrum predictions remain correctly aligned with their
GT counterparts.
Junction-DSC is moderate (0.617) because two of three junction
classes are predicted correctly; the L6 class contributes a near-
zero DSC term (0.0, since no L6 voxels are predicted) which the
3-class average dilutes.

The clinical implication of L6 suppression is more recoverable
than that of level-shift mislabeling: a surgical planner reading
the predicted L4--L5 boundary will target the correct
L4--L5 disc, since L4 and L5 predictions remain correctly
positioned.
The L5--L6 disc, however, is invisible to the planner---the
patient's actual lumbarization is silently absent from the model
output.
A planner asked to target the L5--S1 disc will identify the
predicted L5--sacrum boundary, which is anatomically the
L6--sacrum boundary, again one level off-target.

\subsection{Case 3: Count-Style Sacralization Producing Inverse
            Level Shift (Token from sacr\_count subgroup)}
\label{sec:case_sacr_count}

Count-style sacralization cases ($n=9$, mean L5 DSC = 0.000)
show the inverse failure pattern.
GT here labels four lumbar bodies (L1--L4) plus a fused
L5/sacrum mass labeled as sacrum.
TS predicts five lumbar bodies (L1--L5) in their standard
positions.
The result is a uniform $+1$-level offset: every TS-predicted
lumbar label sits on the GT body of the next-rostral identity.
Per-class DSC collapses across all lumbar classes
(0.265, 0.100, 0.055, 0.001, 0.000).
Junction-DSC remains at 0.697 because the sacrum prediction
(which now extends into what TS labels as ``L5'' but GT labels as
sacrum) overlaps reasonably with the GT sacrum within the 40~mm
window.

This failure mode is interesting because it is symmetric to the
lumbarization case: in lumbarization, GT has $n+1$ lumbar bodies
and TS has $n$ labels; in count-style sacralization, GT has $n-1$
lumbar bodies and TS still has $n$ labels.
Both produce a one-level offset between prediction and ground
truth, but in opposite directions.
The shared root cause is TS's fixed label cardinality.

\subsection{Synthesis}

The three case studies establish a unified picture of TS's
behavior under LSTV.
Architecturally, TS has no mechanism to express the
``six lumbar bodies'' or ``four lumbar bodies + sacralized L5''
hypothesis; its output head is locked to the five-lumbar-bodies
template.
Empirically, the trained model handles violations of the template
in three observable ways: (a) level-shift the entire lumbar label
assignment (Case 1; most common on lumbarization); (b) suppress
the supernumerary body to background (Case 2; less common but
clinically more recoverable); (c) emit standard L1--L5 predictions
that systematically overshoot GT (Case 3; near-universal on
count-style sacralization).
None of these is correct.
All are silent at inference time.

The level-shift pattern (Case 1, Case 3) is the failure mode of
greatest clinical concern: the output is locally well-formed, the
mislabeling is silent, and a downstream planning system has no
mechanism to detect that a one-level shift has occurred.
A counting-aware audit step at the L5/S1 junction---verifying
that the predicted L5--sacrum boundary aligns with the
patient's actual L5--sacrum boundary---would catch this failure,
but no production CT segmentation system to our knowledge
performs such an audit.

% =============================================================================
\section{\dataset{}: Construction and Composition}
\label{sec:dataset}

\paragraph{Patient cohort and source datasets.}
\dataset{} fuses CT scans from the COLONOG
collection~\citep{Smith2015DataCOLONOGRAPHY} with annotations from
CTSpine1K~\citep{Deng2021CTSpine1K:Tomography} (per-vertebra masks via VerSe
labels) and CTPelvic1K~\citep{Liu2021DeepModels} (sacrum and bilateral
hips).
The three resources share the same patient pool but were curated
independently with no SeriesInstanceUID linkage and using
different DICOM-to-NIfTI toolchains, producing silently
incompatible coordinate conventions and partial annotation overlap
(Table~\ref{tab:scan_stats}).
A critical implication is that of 1{,}153 matched records, only
342 carry both spine and pelvis ground truth in fused mode;
352 are separate mode (spine and pelvic from different
acquisitions), 90 are spine-only, and 20 are pelvic-only.

\begin{table}[t]
\caption{\dataset{} composition. Restricting training to fused
cases would discard 70\% of the data, disproportionately
penalizing rare LSTV subgroups; \S\ref{sec:partial} describes the
ignore-label protocol that incorporates the
partially-annotated remainder.}
\label{tab:scan_stats}
\centering
\small
\begin{tabular}{lcc}
\toprule
\bfseries Property & \bfseries Value & \bfseries LSTV cases \\
\midrule
Patients matched (\dataset{})       & 802     & --- \\
Records matched (\dataset{})        & 1{,}153 & 33 patients (44 records) \\
\quad Fused (spine + pelvis)        & 342     & --- \\
\quad Separate (spine + pelvic)     & 352     & --- \\
\quad Spine-only                    & 90      & --- \\
\quad Pelvic-only                   & 20      & --- \\
\midrule
Lumbarization                       & ---     & 14 \\
Sacralization (count-style)         & ---     & 9 \\
Sacralization (full)                & ---     & 6 \\
Semi-sacralization                  & ---     & 2 \\
Ambiguous (cross-anatomy disagreement) & --- & 2 \\
\bottomrule
\end{tabular}
\end{table}

\paragraph{Patient-anchored fusion.}
We resolve patient identity via canonical UID matching (stripping
zero-padding from the final UID component reduces the problem to
an exact relational join, resolving 825/825 patients).
Within each patient's bounded set of TCIA series, we assign each
mask to its source CT acquisition by exhaustive intrapatient
bone-HU maximization:
\begin{equation}
s^*(m) = \argmax_{s \in \mathcal{S}_p}
\frac{|\{v \in \mathcal{L}(m, s) : I_s(v) > 200~\text{HU}\}|}{|\mathcal{L}(m, s)|},
\label{eq:bonehu}
\end{equation}
which serves as a metadata-free oracle since
\texttt{SeriesDescription} and \texttt{BodyPartExamined} fields
are unreliable in TCIA releases.
This achieves 99.6\% spine and 100\% pelvic placement.

\paragraph{Coordinate-convention correction.}
CTSpine1K and CTPelvic1K were converted from DICOM with different
\texttt{dcm2niix} invocations, producing a systematic Y-axis sign
reflection in the resulting NIfTI affines.
For the same patient, CTSpine1K masks place the body center near
$y \approx -181$~mm world coordinate while CTPelvic1K places it
at $y \approx +181$~mm.
This is undetectable from voxel spacing alone and would silently
corrupt all fused masks. We detect via the affine-derived center
sign and reverse the mask along axis~1 with affine update; the
correction is exact and verified on all 714 pelvic masks.
This is an instance of a broad failure class: any pipeline
combining independently processed NIfTI datasets must validate
affine-derived world coordinates before training.

\paragraph{Unified 10-class label map.}
Both source label spaces are remapped:
0\,=\,background, 1--5\,=\,L1--L5,
\textbf{6\,=\,L6 (lumbarized S1)}, 7\,=\,sacrum,
8\,=\,left hip, 9\,=\,right hip.
Class~6 is the first evaluation class in any public CT spine
benchmark reserved for the lumbarized transitional body.
At voxel-level conflicts at the L5/sacrum boundary, CTPelvic1K
takes precedence for sacral extent.

\paragraph{Six-way LSTV taxonomy with dual-origin labels.}
We retain both LSTV label sources separately and partition cases
into six subtypes by combining filename-qualifier evidence
(CTPelvic1K morphological criterion) with vertebral-counting
evidence (CTSpine1K VerSe-label-presence heuristic):
\textbf{normal} (no LSTV indication from either source);
\textbf{lumb} (6 vertebrae labeled or VerSe label 25 present);
\textbf{sacr\_count} (4 vertebrae labeled with sacralization
flag);
\textbf{semisacralization} (filename qualifier
\texttt{\_semisacralization\_});
\textbf{sacralization} (filename qualifier
\texttt{\_sacralization\_}); and \textbf{ambiguous} (cross-
source disagreement that cannot be reconciled by either rule).
Earlier dataset extractions collapsed semisacralization into full
sacralization due to a substring-match parser bug in the qualifier
parser; we identify and fix this, recovering the semi-
sacralization subtype as a distinct evaluation class.
The 6-way taxonomy is essential for the failure-mode analysis in
\S\ref{sec:level_shift}: full sacralization, semi-sacralization,
and count-style sacralization produce qualitatively different
failure patterns and would be invisible if collapsed into a single
``sacralization'' bucket.
We release all labels with origin flags so downstream training
and evaluation can exploit rather than discard the source-level
disagreement.

\subsection{Partial-Annotation Training Protocol}
\label{sec:partial}

Training only on fused cases discards 70\% of the data.
We use nnU-Net v2's
\texttt{ignore\_label}~\citep{Isensee2024Nnu-netSegmentation}
mechanism to train on all 1{,}153 records:
for spine-only and pelvic-only cases, voxels with value 0 are
rewritten to 255 (declared as \texttt{ignore\_label} in
\texttt{dataset.json}); voxels with value 255 contribute zero to
both Dice and cross-entropy loss, so the model is supervised only
on structures present in each case's ground truth.
Spine classes (1--6) receive supervision from all 784
fused-or-spine-only-or-spine-side records;
pelvic classes (7--9) from all 714
fused-or-pelvic-only-or-pelvic-side records.
This is a $3.36\times$ effective training-set expansion relative
to fused-only training.

% =============================================================================
\section{Evaluation Protocol}
\label{sec:protocol}

We define three metrics as the minimum required to attribute LSTV
failure modes:

\paragraph{Per-class DSC stratified by 6-way LSTV subgroup.}
Reported separately for normal, lumb, sacr\_count, semi-
sacralization, sacralization, and ambiguous, on each of 9
anatomical classes plus mean foreground.
This exposes which specific classes degrade on which
phenotypes; it is what \S\ref{sec:level_shift} relied on to
isolate the caudally-progressive lumbar collapse on lumbarization
cases.

\paragraph{Junction-DSC over a 40~mm L5/S1 axial window.}
DSC restricted to the $z$-axis window centered at the L5/S1
boundary (or its lumbarized equivalent in lumbarization cases).
This window is the clinical window of interest---surgical planning
around L5/S1 cares about boundary delineation, not L1 voxel-level
overlap.
Junction-DSC is what reveals the LSTV gap that aggregate per-class
metrics smooth over.

\paragraph{Per-class voxel confusion matrices at the L5/S1
junction.}
For each LSTV-positive case, the ground-truth $\times$ prediction
confusion matrix over the 40~mm window.
This is what \S\ref{sec:case_studies} relies on to expose the
level-shift mechanism: the matrices reveal that the lumbarization
failure is not a diffuse degradation but a structured one-level
offset, which DSC alone cannot characterize.
A model that produces the same junction-DSC by emitting noise
versus by producing a clean one-level shift is qualitatively
different in clinical risk; the confusion matrix is the only
metric that distinguishes them.

% =============================================================================
\section{Full-Resolution Baseline}
\label{sec:baseline}

\paragraph{Architecture and configuration.}
nnU-Net v2 with the ResEnc-M planner, configuration
\texttt{3d\_fullres}, plans
\texttt{nnUNetResEncUNetPlans\_100G} (100~GB GPU-memory target).
The planner selects:
voxel spacing $0.80\!\times\!0.75\!\times\!0.75$~mm,
patch size $\boldsymbol{256\!\times\!320\!\times\!320}$,
batch size 2, 7-stage residual encoder
($n_\text{features}=[32,64,128,256,320,320,320]$;
$n_\text{blocks}=[1,3,4,6,6,6,6]$),
deep supervision enabled,
loss $=$ class-reweighted \texttt{DC\_and\_CE\_loss} (L6 weight
$=$ 4.0, background $=$ 0.5, others $=$ 1.0),
500 epochs, SGD with Nesterov momentum,
PolyLR schedule starting at $10^{-2}$.

\paragraph{LSTV oversampling and L6 patch bias.}
The 4--5\% LSTV base rate is too low for the model to learn the
L6 class robustly without intervention.
We apply two compounding mechanisms.
First, queue-level oversampling replicates lumb, sacr\_count,
and full-sacralization cases until they constitute \textbf{25\%}
of training samples per fold (semi-sacralization and ambiguous
cases are excluded from oversampling due to small $n$ to avoid
memorization, but remain tracked for per-subgroup validation).
Second, foreground-class patch-sampling bias: when nnU-Net's
foreground oversampler samples a class for a lumbarization case,
the L6 class is selected with probability 0.6 (versus
$1/n_\text{classes}\,{\approx}\,0.11$ default), increasing
gradient signal on the class TS structurally cannot learn.
Per-LSTV-subgroup validation Dice is logged during training so
the headline metric (L6 Dice on lumbarization cases) is monitored
continuously across all 6 subgroups.

\paragraph{Hardware.}
H200 NVL (140~GB HBM3e), 24 CPU workers for data augmentation,
nnU-Net's preunpacking step run once at job start with a per-fold
shared local-NVMe cache (1.7~TB \texttt{/tmp}, $\sim$3 GB/s
sequential read; replaces NFS reads which throttled to 21 MB/s
under 5-fold parallel contention).
Per-epoch wall-clock: 350--950~s under contention from 5 parallel
folds; ${\approx}\,400$~s steady-state. Total per-fold wall-clock:
${\approx}\,3.7$ days for 500 epochs.

\paragraph{Why not smaller GPUs?}
The $256\!\times\!320\!\times\!320$ patch is what
\texttt{nnUNetResEncUNetPlans\_100G} selects to maximize memory
utilization at full resolution while fitting both the lumbar spine
and pelvic context.
Reducing GPU memory budget forces either a smaller patch (loses the
spinopelvic spatial relationship the model needs to disambiguate
L5/S1) or coarser spacing (loses the sub-millimeter boundary
detail required for Castellvi Type~II).
Neither tradeoff is acceptable for an LSTV-targeted model.
This is the empirical reason TS-class CT segmentation models
trained at 1.5--3~mm cannot recover the LSTV gap by data
augmentation alone---they lack the resolution to see what they
need to learn.

\paragraph{Splits.}
5-fold patient-stratified CV with 6-way LSTV stratification
(LSTV-first stratification ensures each fold contains a
representative mix of normal, lumb, sacr\_count, sacralization,
semi-sacralization, and ambiguous cases).
Patient-level (not record-level) split assignment guarantees no
data leakage between train and validation across separate-mode
patients.
For the n=2 rare classes (semi-sacralization, ambiguous), only
folds 0--1 contain validation cases; these subgroups are reported
as aggregate-across-folds metrics rather than per-fold means.
Splits and metadata released alongside the dataset.

\subsection{Baseline Results}

\textit{Result tables to be filled from completed 5-fold CV
(ETA: training completes May 5, 2026).}

\begin{table}[t]
\caption{\dataset{} 5-fold CV baseline on the held-out validation
folds, 6-way taxonomy.
$\Delta$\,= absolute improvement over \TS{} zero-shot
(Table~\ref{tab:ts_stratified}).
The baseline closes the LSTV junction-DSC gap and emits an L6
class on all lumbarization cases, confirming that the TS failure
mode is closable when (a) the label vocabulary includes L6 and
(b) training data is LSTV-stratified at full resolution.}
\label{tab:baseline}
\centering
\small
\begin{tabular}{lrcccccr}
\toprule
\bfseries Subgroup & \bfseries $n$ &
\bfseries L1--L5 & \bfseries L6 & \bfseries Sac &
\bfseries Jxn DSC & \bfseries Lumb level shift & \bfseries L6 emitted \\
\midrule
Normal                  & XX & 0.XXX & ---   & 0.XXX & 0.XXX & 0\%   & 0/N \\
Lumbarization           & 14 & 0.XXX & 0.XXX & 0.XXX & 0.XXX & X\%   & XX/14 \\
Sacralization (count)   &  9 & 0.XXX & ---   & 0.XXX & 0.XXX & X\%   & 0/9 \\
Semi-sacralization      &  2 & 0.XXX & ---   & 0.XXX & 0.XXX & 0\%   & 0/2 \\
Sacralization (full)    &  6 & 0.XXX & ---   & 0.XXX & 0.XXX & 0\%   & 0/6 \\
Ambiguous               &  2 & 0.XXX & X.XX  & 0.XXX & 0.XXX & X\%   & X/2 \\
\midrule
$\Delta$ vs.\ TS (lumb)    & --- & +XX.X & +XXX & +X.X & \textbf{+XX.X} & --- & --- \\
$\Delta$ vs.\ TS (sacr\_count) & --- & +XX.X & --- & +X.X & \textbf{+XX.X} & --- & --- \\
\bottomrule
\end{tabular}
\end{table}

% =============================================================================
\section{Discussion}
\label{sec:discussion}

\paragraph{What the protocol reveals.}
The combination of 6-way-stratified per-class DSC, junction-DSC,
and per-class voxel confusion matrices makes architectural and
distributional failure modes individually attributable.
Without all three, the lumbarization failure of TS appears as a
diffuse degradation that could plausibly be fixed by more training
data; with all three, it is unambiguously a combination of
(a) missing L6 class
(\textsc{TS\_L6} structurally zero across all 1{,}152 evaluation
cases) and (b) learned level-shift labeling
(GT-L6 voxels predicted as L5 in 95\%+ of the L6 region on
lumbarization cases; per-class L4--L5 DSC collapse to 0.39 / 0.29).
The first is fixed by changing the label vocabulary; the second
requires LSTV-stratified training data with L6-targeted patch
sampling and sufficient resolution to learn the distinguishing
morphological features.
Our baseline does both and recovers the gap.

\paragraph{The 6-way taxonomy matters.}
Earlier extractions of LSTV labels from CTPelvic1K filename
qualifiers used substring matching, which collapsed the
\texttt{\_semisacralization\_} qualifier into the
\texttt{\_sacralization\_} bucket.
We identify this parser bug and recover semi-sacralization as a
distinct subtype.
The recovery is consequential for the failure-mode analysis: full
sacralization shows a moderate failure (L1--L5 DSCs 0.67--0.79,
junction-DSC 0.636) while semi-sacralization shows essentially no
failure (L1--L5 DSCs $>0.95$).
Collapsing them would have produced a population mean that
obscures both signals and licensed a misleading conclusion that
``sacralization is moderate-difficulty for TS.''
The lesson is that LSTV evaluation cannot use binary
classification (LSTV positive/negative); the failure modes vary
by morphology in ways that affect the prescription.

\paragraph{Generalization beyond TS.}
The level-shift mislabeling failure is a property of any spine
segmentation model whose output head excludes L6, deployed on a
distribution that contains lumbarization.
We expect the same failure mode in MRI spine
models~\citep{Warszawer2025TotalSpineSeg:MRI,GraafSPIDERBenchmark}
that share TS's label vocabulary;
direct evaluation requires modality-converted comparison and is
beyond the scope of the present work.

\paragraph{Hardware as a benchmark contribution.}
Reporting that an evaluation gap is closable is materially
different from reporting that it is closable \emph{at scale that
runs on commodity hardware}.
We document explicitly that the 100~GB patch size is the smallest
that fits the necessary spinopelvic context at the resolution
required to learn Castellvi-defining features, and that this
budget is presently available only on H200-class hardware.
A reader who cannot replicate the 100~GB patch on their hardware
should expect a quantifiably worse LSTV gap closure;
this is not a property of nnU-Net or our specific training
pipeline but of the anatomy.

\paragraph{Limitations.}
COLONOG is a colonoscopy-screening cohort (US, ages 45--75,
academic centers, soft-tissue kernel); generalization to emergency
or pediatric populations is not evaluated.
LSTV labels are inherited from CTPelvic1K and CTSpine1K
heuristics rather than re-read by an independent radiologist---
the cross-source disagreement (full sacralization morphological
qualifier versus count-style label inferred from missing L5
annotation) is a real signal of definitional indeterminacy, but
it also bounds the ground-truth quality of the sacralization
subgroups; we recommend evaluation focus on lumbarization, where
agreement is well-defined.
Two of the six subtypes (semi-sacralization, ambiguous) have
$n=2$ patients each, sufficient for case-level documentation
(\S\ref{sec:case_studies}) but underpowered for population-level
performance estimates; we report them as descriptive case-level
findings rather than statistical claims.
TS was evaluated zero-shot, which is the correct comparison for
``TS as production system'' but undersells what TS could achieve
if retrained with an L6 class on LSTV-stratified data;
we view that retraining as a natural follow-up enabled by
\dataset{}.

\paragraph{Implications.}
A production CT segmentation system that emits a confident
level-shifted prediction is more dangerous than one that fails
noisily.
The infrastructure to detect this---a stratified evaluation
benchmark with the per-class confusion matrix metric---did not
exist before this work.
Any future spine segmentation model claiming production readiness
should report per-subgroup DSC and per-class confusion at the
L5/S1 junction on a stratified LSTV cohort.

\bibliographystyle{plainnat}
\bibliography{bib/references}

\end{document}
